\documentclass[border=2pt]{standalone}
\usepackage{polymer-paper-preamble}

% v0.3 L3 snippet library
%% polymer_chain.snippet.tex
%% Inverse vulcanization sulfur-rich linear polymer — ball-and-stick TikZ schematic.
%% Fixed pattern: C₄-S₄-C₄-S₆-C₄-S₄-C₃ (chemistry of sulfur polymer trap sites).
%% License: snippet code MIT-style.
%% Author origin: figure-agent v0.3 — L3 snippet (ball-and-stick unified, reference-grounded).
%%
%% Usage:
%%   %% polymer_chain.snippet.tex
%% Inverse vulcanization sulfur-rich linear polymer — ball-and-stick TikZ schematic.
%% Fixed pattern: C₄-S₄-C₄-S₆-C₄-S₄-C₃ (chemistry of sulfur polymer trap sites).
%% License: snippet code MIT-style.
%% Author origin: figure-agent v0.3 — L3 snippet (ball-and-stick unified, reference-grounded).
%%
%% Usage:
%%   %% polymer_chain.snippet.tex
%% Inverse vulcanization sulfur-rich linear polymer — ball-and-stick TikZ schematic.
%% Fixed pattern: C₄-S₄-C₄-S₆-C₄-S₄-C₃ (chemistry of sulfur polymer trap sites).
%% License: snippet code MIT-style.
%% Author origin: figure-agent v0.3 — L3 snippet (ball-and-stick unified, reference-grounded).
%%
%% Usage:
%%   \input{styles/snippets/polymer_chain.snippet.tex}
%%   \PolymerChainBallStick{x_anchor}{y_anchor}
%%
%% Args:
%%   #1 = anchor_x (TikZ cm, left end position)
%%   #2 = anchor_y (TikZ cm, chain center y)
%%
%% Visuals (Pyun-style ball-and-stick):
%%   - C atom (organic spacer): black filled circle (0.08cm radius, outline 0.28pt)
%%   - S atom: cAmber filled circle (0.10cm radius, outline 0.28pt)
%%   - C-C bond: black line (0.70pt)
%%   - S-S bond: cAmber line (1.0pt)
%%   - bond length: 0.40cm
%%   - Horizontal span ≈ 5.4cm (for 3-chain stacking with moderate gap)

\makeatletter
\newcommand{\PolymerChainBallStick}[2]{%
  \begingroup
  \pgfmathsetmacro{\xstart}{#1}%
  \pgfmathsetmacro{\yc}{#2}%
  \pgfmathsetmacro{\bondlen}{0.40}%
  \pgfmathsetmacro{\radiusC}{0.080}%
  \pgfmathsetmacro{\radiusS}{0.100}%

  % Helper: draw C atom (black ball) at given position
  \newcommand{\DrawC}[1]{%
    \draw[fill=black, line width=0.28pt, draw=black] (#1,\yc) circle (\radiusC);%
  }%
  % Helper: draw S atom (cAmber ball) at given position
  \newcommand{\DrawS}[1]{%
    \draw[fill=cAmber, line width=0.28pt, draw=cAmber!50] (#1,\yc) circle (\radiusS);%
  }%
  % Helper: draw C-C bond (black line)
  \newcommand{\BondCC}[2]{%
    \draw[line width=0.70pt, black] (#1,\yc) -- (#2,\yc);%
  }%
  % Helper: draw S-S bond (cAmber line)
  \newcommand{\BondSS}[2]{%
    \draw[line width=1.0pt, cAmber] (#1,\yc) -- (#2,\yc);%
  }%

  % Organic spacer C₁-C₂-C₃-C₄ (x = 0, 0.4, 0.8, 1.2)
  \DrawC{0.00}%
  \BondCC{0.00}{0.40}%
  \DrawC{0.40}%
  \BondCC{0.40}{0.80}%
  \DrawC{0.80}%
  \BondCC{0.80}{1.20}%
  \DrawC{1.20}%

  % S₁-S₂-S₃-S₄ run (x = 1.6, 2.0, 2.4, 2.8)
  \BondSS{1.20}{1.60}%
  \DrawS{1.60}%
  \BondSS{1.60}{2.00}%
  \DrawS{2.00}%
  \BondSS{2.00}{2.40}%
  \DrawS{2.40}%
  \BondSS{2.40}{2.80}%
  \DrawS{2.80}%

  % Organic spacer C₅-C₆-C₇-C₈ (x = 3.2, 3.6, 4.0, 4.4)
  \BondCC{2.80}{3.20}%
  \DrawC{3.20}%
  \BondCC{3.20}{3.60}%
  \DrawC{3.60}%
  \BondCC{3.60}{4.00}%
  \DrawC{4.00}%
  \BondCC{4.00}{4.40}%
  \DrawC{4.40}%

  % S₅-S₆-S₇-S₈-S₉-S₁₀ run (x = 4.8, 5.2, 5.6, 6.0, 6.4, 6.8) — TRAP SITE
  \BondSS{4.40}{4.80}%
  \DrawS{4.80}%
  \BondSS{4.80}{5.20}%
  \DrawS{5.20}%
  \BondSS{5.20}{5.60}%
  \DrawS{5.60}%
  \BondSS{5.60}{6.00}%
  \DrawS{6.00}%
  \BondSS{6.00}{6.40}%
  \DrawS{6.40}%
  \BondSS{6.40}{6.80}%
  \DrawS{6.80}%

  % Organic spacer C₉-C₁₀-C₁₁-C₁₂ (x = 7.2, 7.6, 8.0, 8.4)
  \BondCC{6.80}{7.20}%
  \DrawC{7.20}%
  \BondCC{7.20}{7.60}%
  \DrawC{7.60}%
  \BondCC{7.60}{8.00}%
  \DrawC{8.00}%
  \BondCC{8.00}{8.40}%
  \DrawC{8.40}%

  % S₁₁-S₁₂-S₁₃-S₁₄ run (x = 8.8, 9.2, 9.6, 10.0)
  \BondSS{8.40}{8.80}%
  \DrawS{8.80}%
  \BondSS{8.80}{9.20}%
  \DrawS{9.20}%
  \BondSS{9.20}{9.60}%
  \DrawS{9.60}%
  \BondSS{9.60}{10.00}%
  \DrawS{10.00}%

  % Organic terminal C₁₃-C₁₄-C₁₅ (x = 10.4, 10.8, 11.2)
  \BondCC{10.00}{10.40}%
  \DrawC{10.40}%
  \BondCC{10.40}{10.80}%
  \DrawC{10.80}%
  \BondCC{10.80}{11.20}%
  \DrawC{11.20}%

  \endgroup
}
\makeatother

%%   \PolymerChainBallStick{x_anchor}{y_anchor}
%%
%% Args:
%%   #1 = anchor_x (TikZ cm, left end position)
%%   #2 = anchor_y (TikZ cm, chain center y)
%%
%% Visuals (Pyun-style ball-and-stick):
%%   - C atom (organic spacer): black filled circle (0.08cm radius, outline 0.28pt)
%%   - S atom: cAmber filled circle (0.10cm radius, outline 0.28pt)
%%   - C-C bond: black line (0.70pt)
%%   - S-S bond: cAmber line (1.0pt)
%%   - bond length: 0.40cm
%%   - Horizontal span ≈ 5.4cm (for 3-chain stacking with moderate gap)

\makeatletter
\newcommand{\PolymerChainBallStick}[2]{%
  \begingroup
  \pgfmathsetmacro{\xstart}{#1}%
  \pgfmathsetmacro{\yc}{#2}%
  \pgfmathsetmacro{\bondlen}{0.40}%
  \pgfmathsetmacro{\radiusC}{0.080}%
  \pgfmathsetmacro{\radiusS}{0.100}%

  % Helper: draw C atom (black ball) at given position
  \newcommand{\DrawC}[1]{%
    \draw[fill=black, line width=0.28pt, draw=black] (#1,\yc) circle (\radiusC);%
  }%
  % Helper: draw S atom (cAmber ball) at given position
  \newcommand{\DrawS}[1]{%
    \draw[fill=cAmber, line width=0.28pt, draw=cAmber!50] (#1,\yc) circle (\radiusS);%
  }%
  % Helper: draw C-C bond (black line)
  \newcommand{\BondCC}[2]{%
    \draw[line width=0.70pt, black] (#1,\yc) -- (#2,\yc);%
  }%
  % Helper: draw S-S bond (cAmber line)
  \newcommand{\BondSS}[2]{%
    \draw[line width=1.0pt, cAmber] (#1,\yc) -- (#2,\yc);%
  }%

  % Organic spacer C₁-C₂-C₃-C₄ (x = 0, 0.4, 0.8, 1.2)
  \DrawC{0.00}%
  \BondCC{0.00}{0.40}%
  \DrawC{0.40}%
  \BondCC{0.40}{0.80}%
  \DrawC{0.80}%
  \BondCC{0.80}{1.20}%
  \DrawC{1.20}%

  % S₁-S₂-S₃-S₄ run (x = 1.6, 2.0, 2.4, 2.8)
  \BondSS{1.20}{1.60}%
  \DrawS{1.60}%
  \BondSS{1.60}{2.00}%
  \DrawS{2.00}%
  \BondSS{2.00}{2.40}%
  \DrawS{2.40}%
  \BondSS{2.40}{2.80}%
  \DrawS{2.80}%

  % Organic spacer C₅-C₆-C₇-C₈ (x = 3.2, 3.6, 4.0, 4.4)
  \BondCC{2.80}{3.20}%
  \DrawC{3.20}%
  \BondCC{3.20}{3.60}%
  \DrawC{3.60}%
  \BondCC{3.60}{4.00}%
  \DrawC{4.00}%
  \BondCC{4.00}{4.40}%
  \DrawC{4.40}%

  % S₅-S₆-S₇-S₈-S₉-S₁₀ run (x = 4.8, 5.2, 5.6, 6.0, 6.4, 6.8) — TRAP SITE
  \BondSS{4.40}{4.80}%
  \DrawS{4.80}%
  \BondSS{4.80}{5.20}%
  \DrawS{5.20}%
  \BondSS{5.20}{5.60}%
  \DrawS{5.60}%
  \BondSS{5.60}{6.00}%
  \DrawS{6.00}%
  \BondSS{6.00}{6.40}%
  \DrawS{6.40}%
  \BondSS{6.40}{6.80}%
  \DrawS{6.80}%

  % Organic spacer C₉-C₁₀-C₁₁-C₁₂ (x = 7.2, 7.6, 8.0, 8.4)
  \BondCC{6.80}{7.20}%
  \DrawC{7.20}%
  \BondCC{7.20}{7.60}%
  \DrawC{7.60}%
  \BondCC{7.60}{8.00}%
  \DrawC{8.00}%
  \BondCC{8.00}{8.40}%
  \DrawC{8.40}%

  % S₁₁-S₁₂-S₁₃-S₁₄ run (x = 8.8, 9.2, 9.6, 10.0)
  \BondSS{8.40}{8.80}%
  \DrawS{8.80}%
  \BondSS{8.80}{9.20}%
  \DrawS{9.20}%
  \BondSS{9.20}{9.60}%
  \DrawS{9.60}%
  \BondSS{9.60}{10.00}%
  \DrawS{10.00}%

  % Organic terminal C₁₃-C₁₄-C₁₅ (x = 10.4, 10.8, 11.2)
  \BondCC{10.00}{10.40}%
  \DrawC{10.40}%
  \BondCC{10.40}{10.80}%
  \DrawC{10.80}%
  \BondCC{10.80}{11.20}%
  \DrawC{11.20}%

  \endgroup
}
\makeatother

%%   \PolymerChainBallStick{x_anchor}{y_anchor}
%%
%% Args:
%%   #1 = anchor_x (TikZ cm, left end position)
%%   #2 = anchor_y (TikZ cm, chain center y)
%%
%% Visuals (Pyun-style ball-and-stick):
%%   - C atom (organic spacer): black filled circle (0.08cm radius, outline 0.28pt)
%%   - S atom: cAmber filled circle (0.10cm radius, outline 0.28pt)
%%   - C-C bond: black line (0.70pt)
%%   - S-S bond: cAmber line (1.0pt)
%%   - bond length: 0.40cm
%%   - Horizontal span ≈ 5.4cm (for 3-chain stacking with moderate gap)

\makeatletter
\newcommand{\PolymerChainBallStick}[2]{%
  \begingroup
  \pgfmathsetmacro{\xstart}{#1}%
  \pgfmathsetmacro{\yc}{#2}%
  \pgfmathsetmacro{\bondlen}{0.40}%
  \pgfmathsetmacro{\radiusC}{0.080}%
  \pgfmathsetmacro{\radiusS}{0.100}%

  % Helper: draw C atom (black ball) at given position
  \newcommand{\DrawC}[1]{%
    \draw[fill=black, line width=0.28pt, draw=black] (#1,\yc) circle (\radiusC);%
  }%
  % Helper: draw S atom (cAmber ball) at given position
  \newcommand{\DrawS}[1]{%
    \draw[fill=cAmber, line width=0.28pt, draw=cAmber!50] (#1,\yc) circle (\radiusS);%
  }%
  % Helper: draw C-C bond (black line)
  \newcommand{\BondCC}[2]{%
    \draw[line width=0.70pt, black] (#1,\yc) -- (#2,\yc);%
  }%
  % Helper: draw S-S bond (cAmber line)
  \newcommand{\BondSS}[2]{%
    \draw[line width=1.0pt, cAmber] (#1,\yc) -- (#2,\yc);%
  }%

  % Organic spacer C₁-C₂-C₃-C₄ (x = 0, 0.4, 0.8, 1.2)
  \DrawC{0.00}%
  \BondCC{0.00}{0.40}%
  \DrawC{0.40}%
  \BondCC{0.40}{0.80}%
  \DrawC{0.80}%
  \BondCC{0.80}{1.20}%
  \DrawC{1.20}%

  % S₁-S₂-S₃-S₄ run (x = 1.6, 2.0, 2.4, 2.8)
  \BondSS{1.20}{1.60}%
  \DrawS{1.60}%
  \BondSS{1.60}{2.00}%
  \DrawS{2.00}%
  \BondSS{2.00}{2.40}%
  \DrawS{2.40}%
  \BondSS{2.40}{2.80}%
  \DrawS{2.80}%

  % Organic spacer C₅-C₆-C₇-C₈ (x = 3.2, 3.6, 4.0, 4.4)
  \BondCC{2.80}{3.20}%
  \DrawC{3.20}%
  \BondCC{3.20}{3.60}%
  \DrawC{3.60}%
  \BondCC{3.60}{4.00}%
  \DrawC{4.00}%
  \BondCC{4.00}{4.40}%
  \DrawC{4.40}%

  % S₅-S₆-S₇-S₈-S₉-S₁₀ run (x = 4.8, 5.2, 5.6, 6.0, 6.4, 6.8) — TRAP SITE
  \BondSS{4.40}{4.80}%
  \DrawS{4.80}%
  \BondSS{4.80}{5.20}%
  \DrawS{5.20}%
  \BondSS{5.20}{5.60}%
  \DrawS{5.60}%
  \BondSS{5.60}{6.00}%
  \DrawS{6.00}%
  \BondSS{6.00}{6.40}%
  \DrawS{6.40}%
  \BondSS{6.40}{6.80}%
  \DrawS{6.80}%

  % Organic spacer C₉-C₁₀-C₁₁-C₁₂ (x = 7.2, 7.6, 8.0, 8.4)
  \BondCC{6.80}{7.20}%
  \DrawC{7.20}%
  \BondCC{7.20}{7.60}%
  \DrawC{7.60}%
  \BondCC{7.60}{8.00}%
  \DrawC{8.00}%
  \BondCC{8.00}{8.40}%
  \DrawC{8.40}%

  % S₁₁-S₁₂-S₁₃-S₁₄ run (x = 8.8, 9.2, 9.6, 10.0)
  \BondSS{8.40}{8.80}%
  \DrawS{8.80}%
  \BondSS{8.80}{9.20}%
  \DrawS{9.20}%
  \BondSS{9.20}{9.60}%
  \DrawS{9.60}%
  \BondSS{9.60}{10.00}%
  \DrawS{10.00}%

  % Organic terminal C₁₃-C₁₄-C₁₅ (x = 10.4, 10.8, 11.2)
  \BondCC{10.00}{10.40}%
  \DrawC{10.40}%
  \BondCC{10.40}{10.80}%
  \DrawC{10.80}%
  \BondCC{10.80}{11.20}%
  \DrawC{11.20}%

  \endgroup
}
\makeatother

%% log_plot.snippet.tex
%% PGFPlots loglogaxis paper-style — A2 ship.
%% License: pgfplots is LPPL 1.3c (TeX Live); style key code MIT-style.
%%
%% This snippet defines NO new macro. It documents the convention for the
%% pre-existing `paper loglog` style key declared in polymer-paper-preamble.sty,
%% so any consumer who \input-s this file is opting in to the convention and
%% can find usage examples here.
%%
%% Usage:
%%   \begin{loglogaxis}[
%%     paper loglog={3.6cm}{2.6cm},
%%     xmin=1e-3, xmax=1e3, ymin=1e-4, ymax=1e1,
%%     xtick={1e-3, 1e0, 1e3},
%%     ytick={1e-4, 1e-1, 1e1},
%%     xlabel={$\log\,t$}, ylabel={$\log\,I$},
%%   ]
%%     \addplot[cBlue, line width=0.92pt, mark=none] coordinates {...};
%%   \end{loglogaxis}
%%
%% Sizing contract (caller responsibility):
%%   `paper loglog={W}{H}` — W and H are the INNER plot area (the data box).
%%   The style includes `scale only axis`, so adding `title=` / `xlabel=` /
%%   `ylabel=` grows the axis bbox OUTWARD (above/below/beside the data box)
%%   without shrinking the plotted area. To fit a 2.75cm x 2.02cm fixture
%%   target, pass {3.6cm}{2.6cm} — the data area will be exactly that.
%%
%% Pitfall (surfaced 2026-05-04 from golden Row 1):
%%   WITHOUT `scale only axis`, two sibling axes that share `paper loglog={W}{H}`
%%   render with DIFFERENT inner plot heights when one has `title=` and the
%%   other does not — title eats from H. The current style key includes the
%%   fix; do NOT remove it. If you hand-roll your own loglog style, copy
%%   the `scale only axis` keyword.
%%
%% Sampling note: pgfplots `samples=N` uses linear sampling, which clusters
%% points at the high end of a log axis. For deterministic curve placement
%% prefer explicit `coordinates {...}` (log-spaced) over `\addplot {f(x)}`.
%% If function evaluation is needed, set `samples=80, samples y=80`
%% and verify visually. Spike G2 documented this confound.
%%
%% Annotation overlay: place text inside the axis with
%%   \node at (axis cs:<x>, <y>) {...};
%% which uses the data coordinate system (NOT the outer tikzpicture's).
%%
%% Smoke fixture: examples/_snippet_smoke/log_plot/log_plot_smoke.tex
%% First production consumer: examples/golden_trap_depth_picture/golden_trap_depth_picture.tex Row 1.
  % docs only; style key in preamble

\begin{document}
% Required source tokens: Mathematical interpretation; Molecular origin.
\begin{tikzpicture}[
  x=1cm,
  y=1cm,
  every node/.style={font=\sffamily\fontsize{7.2}{8.6}\selectfont, text=cGray!95!black},
  rowlabel/.style={font=\sffamily\fontsize{7.8}{9.4}\selectfont, align=center, text=black},
  titleTeal/.style={font=\sffamily\bfseries\fontsize{9.0}{10.8}\selectfont, text=cTeal},
  tinyLabel/.style={font=\sffamily\fontsize{5.2}{6.2}\selectfont},
  tickLabel/.style={font=\sffamily\fontsize{4.3}{5.2}\selectfont, text=cGray!90!black},
  mathBlue/.style={font=\sffamily\fontsize{8.5}{10.2}\selectfont, text=cBlue},
  axis/.style={-{Stealth[length=3.2pt,width=2.4pt]}, cGray, line width=0.42pt},
  evidenceArrow/.style={-{Stealth[length=4pt,width=2.8pt]}, cGray!70, line width=0.48pt},
  sep/.style={cGray!50, line width=0.32pt},
  trapShallow/.style={cAmber, line width=0.72pt},
  trapDeep/.style={cBlue!45!cRed, line width=0.72pt},
  electron/.style={circle, draw=cBlue!75!black, fill=cBlue!45, inner sep=0pt, minimum size=1.35mm, line width=0.34pt},
  bandbox/.style={draw=none, fill=white, inner sep=0.8pt, minimum width=0.42cm, minimum height=0.20cm},
  bandEdgeCB/.style={cBlue!80!black, line width=1.05pt, line cap=round},
  bandEdgeVB/.style={cBlue!45!cRed!90!black, line width=1.05pt, line cap=round},
]

\newcommand{\SmallLobe}[4]{%
  \draw[#3, line width=0.75pt, fill=#3!18]
    (#1,#2) .. controls (#1+0.18,#2+0.03) and (#1+0.23,#2+0.68) ..
    (#1+0.43,#2+0.69) .. controls (#1+0.63,#2+0.69) and (#1+0.65,#2+0.04) ..
    (#1+0.87,#2) -- cycle;
  \node[text=#3, tinyLabel] at (#1+0.43,#2-0.22) {#4};
}

% Canvas
\fill[white] (-0.05,-0.80) rectangle (16.98,9.22);
\draw[sep] (0.15,5.65) -- (11.95,5.65);
\draw[sep] (0.15,3.45) -- (11.95,3.45);

% Row 1 left icon and label
\begin{scope}[shift={(0.25,6.13)}]
  \draw[cGray, line width=0.55pt, fill=cLGray!45] (0.35,1.35) rectangle (1.18,2.28);
  \draw[black, line width=0.36pt] (0.47,1.65) rectangle (1.06,2.18);
  \draw[cBlue, line width=0.55pt] (0.53,2.09) .. controls (0.68,2.02) and (0.69,1.83) .. (1.00,1.78);
  \foreach \x in {0.58,0.78,0.98} {\draw[cGray, fill=cLGray!20] (\x,1.52) circle (0.035);}
  \draw[cGray, line width=0.50pt] (0.43,1.28) -- (0.58,1.28);
  \node[rowlabel] at (0.78,0.98) {Experiment};
\end{scope}

% Row 1 power-law plot — v0.3 A2: PGFPlots loglogaxis with paper loglog style
\begin{scope}[shift={(3.15,6.15)}]
  \begin{loglogaxis}[
    paper loglog={3.6cm}{2.6cm},
    xmin=1e-3, xmax=1e3, ymin=1e-4, ymax=1e1,
    xtick={1e-3, 1e0, 1e3}, ytick={1e-4, 1e-1, 1e1},
    xlabel={$\log\,t$}, ylabel={$\log\,I$},
  ]
    % I(t) = c * t^{-n}, n=0.7. Log-spaced explicit coords (see snippet docs §sampling).
    \addplot[cBlue, line width=0.92pt, mark=none] coordinates {
      (1e-3, 1.0) (1e-2, 0.2) (1e-1, 0.04) (1e0, 8e-3)
      (1e1, 1.6e-3) (1e2, 3.2e-4) (1e3, 6.4e-5)
    };
    % Slope reference: vertical rise + horizontal run
    \addplot[cBlue, dashed, line width=0.45pt, mark=none]
      coordinates {(1e0, 8e-3) (1e0, 1.6e-3) (1e1, 1.6e-3)};
    \node[font=\sffamily\fontsize{7.0}{8.4}\selectfont, text=cBlue, anchor=north east]
      at (axis cs:1e3, 8e0) {$I(t)\propto t^{-n}$};
    \PlotCallout[text=cBlue, anchor=north west]
      {(axis cs:1e0, 1.6e-3)}{(axis cs:1.2e0, 8e-4)}{slope = $-n$}
  \end{loglogaxis}
\end{scope}

% Row 1 Debye reference — v0.3 A2
% shift x widened from 7.30 → 7.80 (2026-05-04): scale-only-axis moved tick
% labels outside W=3.6cm bbox, so 0.55cm gap caused 10^3 (left) vs 10^-4
% (right) collisions in strict mode. Mid-axis arrow target adjusted below.
\begin{scope}[shift={(7.80,6.02)}]
  \begin{loglogaxis}[
    paper loglog={3.6cm}{2.6cm},
    xmin=1e-3, xmax=1e3, ymin=1e-4, ymax=1e1,
    xtick={1e-3, 1e0, 1e3}, ytick={1e-4, 1e-1, 1e1},
    xlabel={$\log\,t$}, ylabel={$\log\,I$},
    title={Discharge (Debye reference)},
  ]
    % Debye decay: exp(-t/tau), tau=30. Pre-evaluated for stable log sampling.
    \addplot[cGray, line width=0.75pt, mark=none] coordinates {
      (1e-3, 0.99) (1e-2, 0.99) (1e-1, 0.99)
      (1e0, 0.97) (3e0, 0.90) (1e1, 0.72)
      (3e1, 0.37) (1e2, 0.036) (3e2, 4.5e-5)
    };
    \addplot[cGray, dashed, line width=0.38pt, mark=none]
      coordinates {(30, 0.37) (30, 1e-4)};
    \node[align=left, text=cGray, fill=white, inner sep=0.65pt,
      font=\sffamily\itshape\fontsize{6.0}{7.2}\selectfont, anchor=west]
      at (axis cs:3e-3, 1e-2) {Debye\\exp($-t/\tau$)};
    \PlotCallout[text=cGray, anchor=south west]
      {(axis cs:30, 1e-3)}{(axis cs:7e1, 8e-4)}{$\tau_d$}
  \end{loglogaxis}
  % Outgoing arrow to Row 2 — shifted to plot right edge after scale-only-axis fix
  \draw[evidenceArrow] (3.60,1.30) -- (6.14,1.30) -- (6.28,0.90);
\end{scope}

% Row 2 label/icon and flow
\begin{scope}[shift={(0.25,3.62)}]
  \draw[black, line width=0.42pt] (0.42,1.00) rectangle (0.97,1.80);
  \node[font=\sffamily\fontsize{8.0}{9.6}\selectfont, align=center, text=black] at (0.69,1.40) {$\Sigma$\\$=\int$};
  \node[rowlabel] at (0.82,0.43) {Mathematical\\interpretation};
  \node[mathBlue] at (3.00,1.24) {$I(t)\propto t^{-n}$};
  \draw[evidenceArrow] (4.13,1.18) -- (4.88,1.18);
  \node[mathBlue] at (5.18,1.18) {$n$};
  \node[align=center, text=cGray, font=\sffamily\itshape\fontsize{7.0}{8.4}\selectfont] at (6.88,1.45) {Debye\\exp($-t/\tau$)};
  \draw[evidenceArrow] (6.88,1.15) -- (6.88,0.62);
  \node[text=cGray] at (6.88,0.36) {$\tau_d$}; % tau_d
  \draw[evidenceArrow] (7.75,1.18) -- (8.55,1.18);
  \node[text=black, font=\sffamily\fontsize{8.2}{9.8}\selectfont] at (8.95,1.18) {$g(E_t)$}; % g(E_t), E_t
  \begin{scope}[shift={(9.45,0.93)}, xscale=0.72, yscale=0.66]
    \SmallLobe{0}{0}{cAmber}{shallow}
  \end{scope}
  \begin{scope}[shift={(10.05,0.78)}, xscale=1.18, yscale=1.24]
    \SmallLobe{0}{0}{cBlue!45!cRed}{deep}
  \end{scope}
  \draw[evidenceArrow] (11.20,1.18) -- (12.38,1.18);
\end{scope}

% Row 3 molecular origin
\begin{scope}[shift={(0.25,0.08)}]
  \draw[cGray, line width=0.70pt] (0.38,2.30) .. controls (0.20,2.56) and (0.55,2.74) .. (0.78,2.58)
    .. controls (0.60,2.28) and (0.88,2.06) .. (1.16,2.12);
  \foreach \x/\y in {0.47/2.08,0.75/2.02,1.00/2.18} {
    \draw[cAmber, fill=white, line width=0.45pt] (\x,\y) circle (0.045);
  }
  \node[rowlabel] at (0.85,1.42) {Molecular\\origin};

  % v0.3: chemfig-based monomer-resolved chains (replaces 3 plot[smooth] + S \foreach)
  % S branch index map for 11-monomer chain at atom_sep=0.30: monomer i lands at
  % x ~ anchor + 0.52*(i-1). Highlight box (4.20-5.85) covers monomers 6-9.
  \PolymerChain{1.65}{2.40}{11}{3,7,8,9}
  \PolymerChain{1.60}{1.78}{11}{2,6,7,8,9}
  \PolymerChain{1.60}{1.08}{11}{3,7,8,9}
  \draw[cGray, dashed, rounded corners=2pt, line width=0.40pt] (4.20,1.47) rectangle (5.85,2.02);
  \node[text=cAmber, font=\sffamily\itshape\fontsize{7.0}{8.4}\selectfont] at (3.80,0.40) {S-rich segments};
  \draw[evidenceArrow] (4.55,0.62) -- (4.92,1.47);
  \draw[evidenceArrow] (5.05,0.70) -- (5.05,1.43);
  \node[align=center, font=\sffamily\itshape\fontsize{5.6}{6.7}\selectfont, text=black, fill=white, inner sep=1.2pt] at (6.18,0.22) {chemical origin\\(electronegativity,\\polarizability of S)};

  \draw[cGray, dashed, rounded corners=2pt, line width=0.42pt] (7.65,1.12) rectangle (9.60,2.62);
  \node[font=\sffamily\itshape\fontsize{6.5}{7.8}\selectfont, text=black] at (8.63,2.92) {localized traps};
  \TrapLevel{8.02}{2.28}{trapShallow}
  \draw[trapShallow] (8.52,2.24) -- (8.86,2.24);
  \node[text=cAmber] at (9.20,2.24) {$\cdots$};
  \TrapLevel{8.28}{1.55}{trapDeep}
  \draw[trapDeep] (8.80,1.51) -- (9.16,1.51);
  \node[text=cBlue!45!cRed] at (9.48,1.51) {$\cdots$};
  \draw[evidenceArrow] (7.10,1.78) -- (7.55,2.18);
  \draw[evidenceArrow] (9.65,2.20) -- (12.38,2.20) -- (12.40,2.38) -- (12.70,2.38);
  \node[align=center, font=\sffamily\itshape\fontsize{5.6}{6.7}\selectfont, text=black, fill=white, inner sep=1.2pt] at (9.88,0.24) {physical origin\\(local potential\\fluctuations)};
  \draw[evidenceArrow] (7.15,0.52) .. controls (7.55,0.76) and (7.75,1.08) .. (7.82,1.12);
  \draw[evidenceArrow] (9.48,0.55) -- (9.40,1.12);
\end{scope}

% Right-side converged trap-depth picture
\begin{scope}[shift={(12.62,1.95)}, xscale=0.90]
  \draw[cTeal, line width=1.0pt] (0.62,-0.10) .. controls (0.22,-0.10) and (0.22,-0.10) .. (0.22,0.28)
    -- (0.22,5.35) .. controls (0.22,5.73) and (0.22,5.73) .. (0.62,5.73);
  \node[titleTeal, anchor=west] at (1.15,5.75) {converged trap-depth picture};
  \draw[axis] (1.25,1.05) -- (1.25,4.95);
  \node[rotate=90, text=black] at (0.92,3.00) {Energy};
  \draw[bandEdgeCB] (1.70,4.85) -- (3.16,4.85);
  \BandBox{3.35}{4.85}{CB}
  \draw[bandEdgeVB] (1.70,1.10) -- (3.16,1.10);
  \BandBox{3.35}{1.10}{VB}
  \draw[cGray!45, dashed, line width=0.38pt] (3.62,0.95) -- (3.62,5.10);
  \node[text=black] at (3.82,3.12) {$E_t$}; % E_t
  \TrapLevel{2.02}{4.25}{trapShallow}
  \TrapLevel{2.45}{4.12}{trapShallow}
  \TrapLevel{2.06}{3.68}{trapShallow}
  \TrapLevel{2.48}{3.60}{trapShallow}
  \TrapLevel{2.18}{2.42}{trapDeep}
  \TrapLevel{2.50}{2.62}{trapDeep}
  \draw[trapDeep] (1.78,2.02) -- (2.08,2.02);
  \draw[trapDeep] (2.42,1.95) -- (2.75,1.95);
  \node[text=cBlue!45!cRed] at (2.95,2.50) {$\cdots$};

  \draw[axis, black] (4.05,1.52) -- (4.05,4.85);
  \draw[axis, black] (4.05,1.52) -- (5.88,1.52);
  \node[text=black, font=\sffamily\fontsize{8.2}{9.8}\selectfont] at (4.98,5.35) {$g(E_t)$}; % g(E_t)
  \BellCurve[draw=cAmber, line width=0.72pt, fill=cAmber!18]{4.05,3.72,4.72,4.40,side}
  \BellCurve[draw=cBlue!45!cRed, line width=0.72pt, fill=cBlue!45!cRed!18]{4.05,1.56,5.78,3.02,side}
  \node[text=cAmber, font=\sffamily\itshape\fontsize{6.7}{8.0}\selectfont, fill=white, inner sep=0.65pt, anchor=west] at (4.82,4.06) {shallow};
  \node[text=cBlue!45!cRed, font=\sffamily\itshape\fontsize{6.7}{8.0}\selectfont, fill=white, inner sep=0.65pt, anchor=west] at (5.32,2.28) {deep};
\end{scope}

\end{tikzpicture}
\end{document}
